\chapter{Autoimmune Diseases}

\section*{Patient focus}
Autoimmune diseases occur when the body's immune system attacks its own tissues. This can affect almost any organ---joints, lungs, kidneys, brain, blood vessels, and more. Imaging helps doctors see where inflammation is occurring, how much damage has been done, and whether treatment is working. Because autoimmune conditions can mimic many other diseases (including cancer and infection), imaging findings must always be interpreted alongside blood tests and symptoms.


\label{ch:ch67_autoimmune}

\section{Imaging}
Autoimmune disease can affect multiple organs and may mimic infection, malignancy, or treatment-related change. Imaging should be interpreted with symptoms, serology, medication exposure, disease duration, and prior studies.

\section{Systemic Lupus Erythematosus (SLE)}

\subsection{CNS Lupus}
\begin{itemize}
  \item \textbf{MRI findings}: T2/FLAIR hyperintense lesions in subcortical white matter, basal ganglia, or brainstem---may mimic MS
  \item \textbf{Differential}: cerebrovascular disease (SLE causes accelerated atherosclerosis), demyelination, infection (immunosuppressed patients), posterior reversible encephalopathy syndrome (PRES) from hypertension or cyclophosphamide
  \item \textbf{MRA}: vasculitis of medium/small vessels---rarely proven on imaging alone
\end{itemize}

\subsection{Pleural and Pericardial Disease}
\begin{itemize}
  \item \textbf{Pleural effusion}: bilateral, small to moderate; may be the presenting feature
  \item \textbf{Pericardial effusion}: most common cardiac manifestation of SLE; may lead to tamponade
  \item \textbf{CT/MRI}: pericardial thickening and enhancement in active serositis
\end{itemize}

\section{Rheumatoid Arthritis}

\subsection{Hands and Wrists}
\begin{itemize}
  \item \textbf{X-ray findings} (in order of appearance):
    \begin{enumerate}
      \item Periarticular osteopenia
      \item Joint space narrowing (uniform---unlike OA which is asymmetric)
      \item Marginal erosions (at bare areas of bone)
      \item Soft tissue swelling
      \item Deformities (ulnar deviation, swan-neck, boutonniere)
    \end{enumerate}
  \item \textbf{Early disease}: MRI is more sensitive---bone edema (pre-erosive), synovial enhancement (active synovitis), tenosynovitis
  \item \textbf{Typical distribution}: MCP joints 2--5, PIP joints, wrists (especially radiocarpal and midcarpal), MTP joints
\end{itemize}

\subsection{Cervical Spine}
\begin{itemize}
  \item \textbf{Atlantoaxial subluxation} (C1-C2): erosion of odontoid process, increased atlantodental interval ($>$3 mm in adults)
  \item \textbf{Vertical subluxation} (``cranial settling''): odontoid projects above McGregor's line---risk of brainstem compression
  \item \textbf{Clinical significance}: pre-operative cervical spine imaging before intubation in RA patients
\end{itemize}

\begin{pearlbox}[RA hands---key differentiators from OA]
RA shows \textbf{symmetric} joint space narrowing, \textbf{marginal} erosions, and periarticular osteopenia. OA shows \textbf{asymmetric} narrowing, \textbf{osteophytes}, and subchondral sclerosis---no erosions. If erosions are at the DIP joints, think erosive OA rather than RA.
\end{pearlbox}

\section{Vasculitis}

\subsection{Large Vessel Vasculitis}
\begin{table}[h]
\centering
\caption{Large vessel vasculitis syndromes}
\begin{tabular}{|p{0.24\textwidth}|p{0.26\textwidth}|p{0.36\textwidth}|}
\hline
\textbf{Disease} & \textbf{Vessels} & \textbf{Key imaging findings} \\
\hline
Giant cell arteritis & Temporal, vertebral, ophthalmic & Wall thickening, halo sign on US; age $>$50 \\
Takayasu arteritis & Aorta and branches & Wall thickening, stenosis, aneurysm; young females \\
\hline
\end{tabular}
\end{table}

\begin{itemize}
  \item \textbf{CT/MR angiography}: concentric wall thickening of aorta and major branches (``macaroni sign''), stenosis, occlusion, aneurysm
  \item \textbf{FDG-PET/CT}: increased FDG uptake in vessel walls indicates active inflammation---useful for monitoring treatment response and identifying skip lesions
  \item \textbf{Temporal artery US}: halo sign (hypoechoic vessel wall thickening)---non-invasive diagnostic test for GCA
\end{itemize}

\subsection{Small Vessel Vasculitis}
\begin{itemize}
  \item \textbf{ANCA-associated vasculitis} (GPA, MPA, EGPA): imaging depends on organ involvement
    \begin{itemize}
      \item Lungs: nodules (GPA---may cavitate), ground-glass (MPA), eosinophilic pneumonia (EGPA)
      \item Sinuses: mucosal thickening, bone erosion (GPA)
      \item Kidneys: not typically imaged; biopsy is diagnostic
    \end{itemize}
  \item \textbf{IgA vasculitis (Henoch-Schonlein)}: bowel wall thickening (``target sign''), mesenteric edema, renal involvement
\end{itemize}

\section{Inflammatory Myopathy (Myositis)}
\begin{itemize}
  \item \textbf{MRI}: T2 hyperintensity in affected muscles (edema pattern), muscle atrophy in chronic disease; most commonly affects proximal muscles (thighs, shoulders)
  \item \textbf{Patterns}:
    \begin{enumerate}
      \item Dermatomyositis: perifascicular edema
      \item Polymyositis: diffuse muscle edema
      \item Inclusion body myositis: asymmetric involvement, forearm flexors, knee extensors
    \end{enumerate}
  \item \textbf{CT}: muscle atrophy and fatty replacement in chronic disease
  \item \textbf{Association}: consider malignancy screening (CT chest/abdomen/pelvis) in dermatomyositis---15--30\% have underlying malignancy
\end{itemize}

\section{IgG4-Related Disease}
A fibroinflammatory condition that can affect virtually any organ:

\begin{table}[h]
\centering
\caption{IgG4-RD imaging patterns by organ}
\begin{tabular}{|p{0.23\textwidth}|p{0.65\textwidth}|}
\hline
\textbf{Organ} & \textbf{Imaging features} \\
\hline
Aorta & Inflammatory aortic aneurysm \\
Biliary tree & Wall thickening (IgG4-sclerosing cholangitis) \\
Kidneys & Peripheral low-attenuation lesions on CT \\
Lungs & Solitary/multiple nodules, bronchial wall thickening \\
Orbit & Lacrimal gland enlargement, orbital masses \\
Pancreas & Diffuse ``sausage-shaped'' enlargement, ``capsule-like'' rim \\
Retroperitoneum & ``Retroperitoneal fibrosis''---soft tissue encasing ureters, aorta \\
\hline
\end{tabular}
\end{table}

\begin{pearlbox}[IgG4-RD mimics cancer]
IgG4-RD can mimic malignancy in nearly every organ---pancreatic cancer, cholangiocarcinoma, retroperitoneal sarcoma. Key distinguishing features: diffuse organ enlargement rather than focal mass, ``capsule-like'' rim enhancement, bilateral or multifocal involvement, and elevated serum IgG4. Biopsy with IgG4 immunostaining is diagnostic when imaging is equivocal.
\end{pearlbox}

\section{Patterns}
Describe distribution, active inflammation, fibrosis, erosion, vascular involvement, serositis, organ enlargement, and complications. MRI, CT, ultrasound, radiography, and PET may each be useful in different organ systems.

\section{When to Recommend Biopsy}
Biopsy is recommended when:
\begin{enumerate}
  \item Imaging cannot distinguish autoimmune disease from malignancy or infection
  \item Atypical imaging features for the suspected autoimmune diagnosis
  \item Need to confirm IgG4-RD for treatment decisions (rituximab vs observation)
  \item Assessing disease activity in ambiguous cases (e.g.\ FDG-PET-positive vessels vs atherosclerosis)
  \item Vasculitis diagnosis (small vessel---renal biopsy; medium vessel---affected organ biopsy)
\end{enumerate}

\section{Monitoring Treatment Response}
\begin{itemize}
  \item \textbf{MRI}: assesses muscle edema resolution in myositis, synovitis resolution in RA, CNS lesion activity in lupus
  \item \textbf{FDG-PET/CT}: monitors large vessel vasculitis activity (FDG uptake decreases with treatment)
  \item \textbf{CT}: monitors organ involvement in IgG4-RD (size reduction with rituximab)
  \item \textbf{X-ray (hands)}: serial films in RA to assess erosion progression---though MRI is more sensitive for early change
\end{itemize}

\section{Clinical Integration}
Do not infer disease activity from one nonspecific finding. Recommendations should identify the next question, such as targeted MRI, tissue sampling, vascular assessment, or interval follow-up, and should be proportionate to risk.

\begin{pearlbox}[Autoimmune imaging pearls]
\begin{enumerate}
  \item Always correlate imaging with serology (ANA, RF, anti-CCP, ANCA, IgG4 levels, ESR/CRP)
  \item Immunosuppressed patients may have blunted imaging findings---a normal scan does not exclude active disease
  \item Drug-related complications (osteonecrosis, infections, PRES) may mimic or coexist with disease activity
  \item Many autoimmune imaging findings are non-specific---multidisciplinary discussion is essential
\end{enumerate}
\end{pearlbox}

\clearpage
